\icode{E-MD-SI-L2-261119-003}
在环绕地球运动的航天器是处于完全失重的状态，不能利用天平称量物体的质量。当力一定时，物体的质量越大，加速度就越小，根据牛顿第二定律能得出物体的质量。如图所示，已知标准物 \(A\) 的质量为 \(m_1\)，连接在标准物 \(A\) 前后的力学传感器的质量均为 \(m_2\)，待测质量的物体 \(B\) 连接在后传感器上，在某一外力作用下整体在空间站内的桌面上运动，稳定后前后传感器的读数分别为 \(F_1\)、\(F_2\)，由此可知待测物体 \(B\) 的质量为\paren[]

\begin{tikzpicture}[>=stealth,font=\scriptsize]
	\draw[fill=blue!30] (1.2,0) rectangle(.6,1.2);
	\draw[fill=red!30] (3,0) rectangle(3.5,1.2);
	\filldraw[pattern=north east lines] (3,0) rectangle(2.8,1.2) (3.5,0) rectangle(3.7,1.2);
	\draw[purple] (1.2,.6)--(2.8,.6);
	\draw[-stealth] (3.7,.6)--++(.8,0);
	\draw[gray] (2.9,.1)--++(-120:.5)node[left]{传感器} (3.6,.1)--++(-60:.5)node[right]{传感器};
	\node at(.9,.6){\(B\)};
	\node at(3.25,.6){\(A\)};
\end{tikzpicture}
\begin{choices}
	\item \(\dfrac{F_2(m_1 + 2m_2)}{F_1}\)
	\item \(\dfrac{F_1(m_1 + 2m_2)}{F_2}\)
	\item \(\dfrac{F_2(m_1 + 2m_2)}{F_1 - F_2}\)
    \item \(\dfrac{F_1(m_1 + 2m_2)}{F_1 - F_2}\)
\end{choices}
\begin{answer}
	C\qquad 提示：
\end{answer}